\documentclass[11pt,a4paper]{ltjsarticle}
\usepackage{amsmath,amssymb}
\usepackage{graphicx}
\usepackage{float}
\usepackage{booktabs}
\usepackage{array}
\usepackage{url}
\usepackage[left=22mm,right=22mm,top=25mm,bottom=25mm]{geometry}
\usepackage{luatexja-fontspec}
\setmainjfont{HaranoAjiMincho-Regular}
\setsansjfont{HaranoAjiGothic-Medium}
\linespread{1.05}\selectfont

% Recommended PDF filename:
% 2611193_HeizoNakai_visual_media_processing1_dai6kai.pdf
% Example build command:
% latexmk -g -lualatex -jobname=2611193_HeizoNakai_visual_media_processing1_dai6kai main.tex

\begin{document}

\begin{center}
{\LARGE \textbf{視覚メディア処理I 第6回レポート}}
\end{center}
\vspace{1em}

\begin{flushright}
提出日: \today \\
学籍番号: 2611193 \\
氏名: 中井 平蔵
\end{flushright}

\subsection*{1. ARマーカーに求められる設計仕様}

\subsubsection*{1.1 画像中から検出できること}

画像中から検出できなければ識別や姿勢推定に進めないため、ARマーカーは背景から検知しやすい形状を持つ必要がある。
白黒のコントラストが大きい外枠を設ければ、複雑な背景でも輪郭を検出しやすい。

\subsubsection*{1.2 ARマーカーのIDを識別できること}

ARマーカーごとに異なるARコンテンツを対応付けるため、複数のARマーカーを使う場合はそれぞれを区別できる情報が必要である。
内部の点配置、白黒パターン、文字や記号を使うことで、各ARマーカーに異なるIDを割り当てられる。

\subsubsection*{1.3 ARマーカーの向きを判別できること}

画像中で見えているマーカーの向きと実際の上下左右を対応付けるため、ARマーカーは回転しても方向を一意に判断できる必要がある。
そのため、内部パターンは回転対称にしないことが重要である。

\subsubsection*{1.4 カメラ姿勢推定に利用できる特徴点を持つこと}

カメラ姿勢を推定するため、ARマーカー上の位置が分かっている点と、画像上で検出された点を対応付ける必要がある。
四角形では四隅、円形では円周や中心、内部パターンの点などが使われる。

\subsection*{2. 各ARマーカー方式}

\subsubsection*{2.1 四角形と内部パターンを組み合わせたARマーカー}

\begin{figure}[H]
\centering
\includegraphics[width=0.68\linewidth]{/Users/nakaiheizou/Desktop/スクリーンショット 2026-06-29 17.20.32.png}
\caption{四角形の外枠と内部パターンを組み合わせたARマーカーの例}
\label{fig:square-pattern-marker}
\end{figure}

\begin{itemize}
\item 画像中からの検出: 外枠により四角形を検出しやすい。
\item IDの識別: 内部パターンの違いで区別。
\item 向きの判別: 非対称な配置にすれば回転方向を判別。
\item カメラ姿勢推定: 四隅を対応点として使えるため精度が高い。
\end{itemize}

\subsubsection*{2.2 四角形と文字・記号を組み合わせたARマーカー}

\begin{figure}[H]
\centering
\includegraphics[width=0.36\linewidth]{/Users/nakaiheizou/Desktop/スクリーンショット 2026-06-29 17.20.38.png}
\caption{四角形の枠内に文字を配置したARマーカーの例}
\label{fig:square-text-marker}
\end{figure}

\begin{itemize}
\item 画像中からの検出: 外枠により四角形を検出しやすい。
\item IDの識別: 文字や記号そのものをIDとして使える。
\item 向きの判別: 文字の非対称性を利用できる。
\item カメラ姿勢推定: 四隅を使えるが、ID認識は文字の見え方に左右される。
\end{itemize}

\subsubsection*{2.3 円形と内部パターンを組み合わせたARマーカー}

\begin{figure}[H]
\centering
\includegraphics[width=0.80\linewidth]{/Users/nakaiheizou/Desktop/スクリーンショット 2026-06-29 17.20.52.png}
\caption{円形の外周と内部パターンを組み合わせたARマーカーの例}
\label{fig:circle-marker}
\end{figure}

\begin{itemize}
\item 画像中からの検出: 円形外周により領域を検知しやすい。
\item IDの識別: 内部パターンで区別する。
\item 向きの判別: 外周だけでは不十分で、内部の非対称性が必要である。
\item カメラ姿勢推定: 中心や楕円形状は使えるが、四角形ほど明確な対応点を取りにくい。
\end{itemize}

\subsection*{3. 各方式の比較}

\begin{itemize}
\item 四角形＋内部パターン
  \begin{itemize}
  \item 長所: 四隅を明確な対応点として利用できるため、高精度な姿勢推定が可能である。また、主に四角形の輪郭を検出するため、計算量を抑えやすく、高速な処理に適している。平面に印刷して貼り付けるだけで使用でき、設置も容易である。
  \item 短所: 外枠や角の一部が遮蔽されると、検出精度が低下しやすい。また、遠距離からの撮影や、マーカーがカメラに対して大きく傾いた場合には、形状が歪むため認識性能が低下する。
  \end{itemize}
\item 四角形＋文字・記号
  \begin{itemize}
  \item 長所: 四角形の外枠によってマーカー領域を安定して検出できる。さらに、文字や記号をIDとして利用できるため、機械だけでなく人間にとっても識別しやすい。
  \item 短所: 文字や記号を認識する処理が必要となるため、内部パターンのみを用いる方式と比べて計算量が増加する。また、低解像度、画像のぼけ、照明変化などの影響を受けやすく、互いに識別しやすい文字や記号を選定する必要がある。
  \end{itemize}
\item 円形＋内部パターン
  \begin{itemize}
  \item 長所: 円形は回転による外形の変化が小さいため、マーカーの回転に対して比較的頑健である。また、多少の画像のぼけにも対応しやすく、円形物体や回転体への設置にも適している。
  \item 短所: 四角形の四隅のような明確な対応点を得にくいため、姿勢推定には追加の工夫が必要である。また、撮影時には円が楕円として観測される場合があり、円または楕円を検出する処理が必要となる。
  \end{itemize}
\end{itemize}

\end{document}
